\section{Controlled tasks and programmable slots}\label{sec:slots}

We first construct circuits whose one-qubit gates can be selected from arbitrary
tables. This separates the geometry of the circuit from the information about
the target unitary. Section~\ref{sec:compilation} will store finite approximations
to these tables in one Boolean oracle.

All vector norms are Euclidean, and $\norm{A}$ denotes the operator norm of a
linear map $A$. Both $A^*$ and $A^\dagger$ denote its adjoint. We never minimize
this norm over scalar phases. For
$1\le d\le 2^n$, write
\[
 \mathcal H_d=\operatorname{span}\{\ket{0},\ldots,\ket{d-1}\}
 \subseteq (\C^2)^{\otimes n}.
\]
The recursion must synthesize many unitaries coherently. Its input therefore
includes a \emph{label register} with basis $\ket{\lambda}$,
$\lambda\in\{0,1\}^t$, and its prescribed action is
\[
 \mathcal U=\sum_{\lambda}\ket{\lambda}\bra{\lambda}\otimes U_\lambda,
 \qquad U_\lambda\colon\mathcal H_d\longrightarrow\mathcal H_d
 \text{ unitary}.
\]
Let $\iota$ append all working qubits in state $\ket{0}$. A circuit $C$
solves this controlled task with error $\delta$ if
\begin{equation}\label{slot:clean-error}
 \norm{C\iota-\iota\mathcal U}\le\delta.
\end{equation}
Thus the estimate includes restoration of the working qubits, and is uniform
over superpositions of labels. No behavior is prescribed outside the input
subspace $\operatorname{span}\{\ket{\lambda}\}\otimes\mathcal H_d$.
The same estimate holds after tensoring with the identity on any further
register, because operator norm is unchanged by that operation.

A \emph{programmable slot} specifies one target qubit and an ordered list of
other qubits as controls. Given a table of matrices $W_y\in U(2)$, its action is
\begin{equation}\label{slot:definition}
 \sum_{y\in\{0,1\}^c}\ket{y}\bra{y}\otimes W_y,
\end{equation}
where $c$ is the number of controls. A \emph{circuit skeleton} specifies the
registers, targets, and control lists of its slots, but leaves their tables
unspecified. The tables may depend on the target unitaries. Their contents
are not part of the skeleton that the uniform generator must produce.
Reversing a skeleton reverses its slot order and replaces every table by its
adjoint table.
Slots are uniformly controlled one-qubit gates, or quantum multiplexors,
in the terminology of~\cite{BVMS2005}.

Every diagonal unitary on $u\ge1$ qubits uses a single slot: choose one of
the qubits as target and put the two diagonal entries for each value of the
other $u-1$ bits in the corresponding $2\times2$ diagonal table entry.
Additional labels can be included among the controls. In particular,
reflections in subspaces spanned by specified computational basis states
take one slot. A phase depending only on labels can be applied to a zero
working qubit, or by scalar matrices on an existing target qubit.

\begin{lemma}[State preparation]\label{slot:preparation}
For every $u\ge1$ there is a skeleton of $u+1$ slots that maps
$\ket{\lambda}\ket{0^u}$ to $\ket{\lambda}\ket{a_\lambda}$ for any prescribed
family of unit vectors $a_\lambda\in\C^{2^u}$. Its structure depends only on
the register sizes.
\end{lemma}

\begin{proof}
Fix a label and write $a=\sum_x a_x\ket{x}$. For a binary prefix $p$, let
$P(p)=\sum_{x\text{ extending }p}|a_x|^2$. Process the bits from first to last.
At the next bit, controlled on the preceding prefix $p$, apply a real
rotation whose first column is
\[
 \begin{pmatrix}
  \sqrt{P(p0)/P(p)}\\[2pt]
  \sqrt{P(p1)/P(p)}
 \end{pmatrix}
 \quad\text{when }P(p)>0.
\]
Choose any rotation when $P(p)=0$. After $u$ such slots the amplitude at $x$
is $|a_x|$. One diagonal slot then supplies the phases of the $a_x$, choosing
arbitrary phases at zero entries. Including $\lambda$ among the controls
makes all these choices simultaneously. This is the usual conditional
rotation construction for state preparation; here an entire uniformly
controlled rotation is counted as one slot; see~\cite{MVBS2005,BVMS2005}.
\end{proof}

\begin{lemma}[Basis permutations]\label{slot:permutation}
Any family of permutations of the $2^u$ computational basis states can be
implemented with $5u$ slots and $u$ temporary qubits that start and finish
at zero. The skeleton is independent of the permutations.
\end{lemma}

\begin{proof}
For the permutation $\pi_\lambda$, use $u$ slots to write its image bitwise
into a temporary register:
\[
 \ket{\lambda}\ket{x}\ket{0}
 \longmapsto \ket{\lambda}\ket{x}\ket{\pi_\lambda(x)}.
\]
Each slot chooses either $I$ or the bit flip, controlled on $\lambda$ and
$x$. Next use $u$ slots to xor $\pi_\lambda^{-1}(y)$ into the first register,
controlled on the second register $y$. The two registers become
$\ket{0}\ket{\pi_\lambda(x)}$. Swap them using three CNOTs per pair of bits.
\end{proof}

The same convention makes a Fourier transform inexpensive. We use the
positive-phase transform, so that a frequency basis vector becomes its
character in the computational basis.

\begin{lemma}[Fourier transforms]\label{slot:fourier}
The unitary
\[
 \ket{x}\longmapsto 2^{-b/2}\sum_{z=0}^{2^b-1}
       e^{2\pi i xz/2^b}\ket{z}
\]
on $b\ge1$ qubits has a uniformly generated skeleton with $O(b)$ slots.
Products of such transforms have a number of slots linear in their total
register width.
\end{lemma}

\begin{proof}
Write the input integer as $x$. Process its bits in decreasing order of
place value. At the bit of place value $2^{b-1-j}$, apply $H$, then a
diagonal slot controlled on the still unprocessed lower bits. Give the
state $1$ the additional phase that, together with its sign from $H$,
makes the ratio of its two amplitudes equal to
$e^{2\pi i x/2^{b-j}}$. Bits above the current bit contribute integral
multiples of $2\pi$ to this expression, and are no longer needed as
controls. The resulting bit is the Fourier output bit of place value
$2^j$. A final reversal of bit order gives the stated transform. There are
two slots per processed bit and $O(b)$ CNOTs for the reversal.
\end{proof}

These constructions retain the scalar phase of every label branch. This is
essential: phases that would be irrelevant for one isolated unitary become
relative phases when the same operation is controlled by a quantum register.
